\documentclass[sigconf]{acmart}

%% Rights management
\setcopyright{none}
\acmConference[AI4HPCC '26]{Workshop on AI for High Performance Compiler Construction}{July 6--9, 2026}{Belfast, Northern Ireland, UK}

\usepackage{booktabs}
\usepackage{subcaption}
\usepackage{xspace}
\usepackage{listings}
\usepackage{xcolor}

\lstset{
  basicstyle=\ttfamily\scriptsize,
  keywordstyle=\color{blue}\bfseries,
  commentstyle=\color{gray},
  stringstyle=\color{red},
  breaklines=true,
  frame=single,
  columns=fullflexible,
  keepspaces=true,
  numbers=left,
  numberstyle=\tiny\color{gray},
  xleftmargin=1.5em,
}

\newcommand{\sysname}{TritonPilot\xspace}
\newcommand{\pet}{PET\xspace}

\begin{document}

\title{Analysis-Guided LLM Code Generation for GPU Kernels:\\Bridging Polyhedral Analysis and Triton Synthesis}

\author{Qin Xiao}
\affiliation{
  \institution{University of Leeds}
  \city{Leeds}
  \country{United Kingdom}
}
\email{ppnh4382@leeds.ac.uk}

\author{Chunwei Xia}
\affiliation{
  \institution{University of Leeds}
  \city{Leeds}
  \country{United Kingdom}
}

\author{Zheng Wang}
\affiliation{
  \institution{University of Leeds}
  \city{Leeds}
  \country{United Kingdom}
}

\begin{abstract}
Large language models (LLMs) can generate GPU kernel code, but without
understanding parallelization constraints they frequently produce
incorrect or inefficient implementations.  We present \sysname, a
system that extracts parallelization properties from C loop nests
using polyhedral analysis (\pet) and LLVM's dependence analysis, then
encodes them as structured natural-language guidance in the LLM prompt.
On 30 Polybench/C kernels, \sysname achieves 100\% correctness
(vs.\ 93\% without analysis) with 2.15$\times$ median speedup over
single-threaded C, rising to 28.97$\times$ at 8$\times$ problem sizes.
On 151 TSVC loop kernels, \sysname passes 150/151 (99\%);
at 32\,MB per array, the generated GPU kernels achieve 46$\times$
median speedup over single-threaded C and 9$\times$ over
80-thread OpenMP.
The approach generalizes to 8 application kernels from Rodinia and
ECP proxy applications, achieving up to 25.8$\times$ speedup.
\end{abstract}

\begin{CCSXML}
<ccs2012>
  <concept>
    <concept_id>10011007.10011006.10011041</concept_id>
    <concept_desc>Software and its engineering~Compilers</concept_desc>
    <concept_significance>500</concept_significance>
  </concept>
  <concept>
    <concept_id>10010147.10010169.10010170</concept_id>
    <concept_desc>Computing methodologies~Parallel computing methodologies</concept_desc>
    <concept_significance>300</concept_significance>
  </concept>
</ccs2012>
\end{CCSXML}

\ccsdesc[500]{Software and its engineering~Compilers}
\ccsdesc[300]{Computing methodologies~Parallel computing methodologies}

\keywords{LLM code generation, GPU kernels, Triton, polyhedral analysis}

\maketitle

%% ===================================================================
\section{Introduction}
\label{sec:intro}

GPU programming remains a significant barrier to adopting hardware
accelerators for scientific computing.  Writing efficient GPU kernels
requires expertise in thread-level parallelism, memory coalescing,
and synchronization---skills that domain scientists often lack.
OpenAI's Triton~\cite{tillet2019triton} simplifies GPU programming
with a block-level abstraction, but generating correct Triton code
still requires understanding which loop dimensions are safe to
parallelize, how data dependences constrain execution order, and when
synchronization barriers are needed.

Recent work has shown that LLMs can generate functional GPU
code~\cite{chen2024chatgpt, ouyang2024kernelbench}, but they lack the
ability to perform precise dependence analysis.  As a result,
LLM-generated GPU code frequently suffers from:
\begin{enumerate}
  \item \textbf{Correctness failures}: data races from parallelizing
    dimensions with carried dependences;
  \item \textbf{Performance anti-patterns}: serial launches,
    scalar loops, or excessive kernel invocations;
  \item \textbf{Synchronization errors}: missing memory barriers
    between phases in stencil computations.
\end{enumerate}

We present \sysname, which integrates classical compiler analysis into
the LLM prompting pipeline.  \sysname uses \pet~\cite{verdoolaege2012polyhedral}
and LLVM's dependence analysis~\cite{llvm} to extract parallelization
properties, then translates these into structured guidance that steers
the LLM toward correct Triton implementations.  The key insight is that
\emph{analysis results can be communicated to LLMs as domain-expert
advice}, enabling the model to make informed parallelization decisions
without performing dependence analysis itself.  Our contributions:
\begin{itemize}
  \item A pipeline that extracts parallelization constraints via
    polyhedral and LLVM analysis and encodes them as structured LLM
    prompt guidance (\S\ref{sec:method});
  \item Evaluation on 30 Polybench/C kernels (100\% pass, 2.15$\times$
    median), 151 TSVC kernels (99\% pass, 9$\times$ median over
    80-thread OpenMP at 32\,MB), and 8 application kernels from
    Rodinia/ECP proxy applications (\S\ref{sec:eval}).
\end{itemize}


%% ===================================================================
\section{Method}
\label{sec:method}

Given a C loop nest, \sysname proceeds in four stages: (1)~static
analysis, (2)~prompt construction, (3)~LLM generation and validation,
and (4)~iterative refinement.

\paragraph{Static Analysis.}
\pet{} extracts a polyhedral representation from which we compute
parallelizable dimensions with validity constraints, WAR dependences
with loop-level scoping, scalar expansion opportunities, and reduction
types.  When \pet{} cannot handle a construct (e.g., non-affine
subscripts), we fall back to LLVM's \texttt{DependenceAnalysis} pass,
parsing direction vectors to determine which loop levels carry
dependences.

\paragraph{Prompt Construction.}
Analysis results are translated into structured natural-language
sections appended to the LLM prompt.  Five principal categories of
guidance are generated:

\emph{Parallelization dimensions}: Each loop dimension is classified as
valid (safe to parallelize), invalid with a specific reason, or
reduction-safe.  The prompt presents valid options with grid
configurations and flags unsafe dimensions with explanations:
\begin{lstlisting}[language={},frame=none,numbers=none,xleftmargin=0em]
Option 1: Parallelize dim `j` (innermost)
  - Grid: triton.cdiv(NY, BLOCK)
  - Each CTA processes BLOCK elements of j
  - Issue: None (fully parallel)
Option 2: Parallelize dim `i` (outermost)
  - Issue: Dependencies carried on `i`
\end{lstlisting}

\emph{WAR dependences}: \pet{} computes exact WAR sets via polyhedral
dependence analysis; when \pet{} fails, we parse LLVM direction vectors
and map them to loop levels.  Three cases are distinguished---no WAR
(in-place update), WAR with parallel resolution (\texttt{B = A.clone()}
and read from clone while writing to original), and WAR requiring
sequential execution (\texttt{grid=(1,)} with explicit loops inside
the kernel).

\emph{Multi-phase decomposition}: Many kernels contain multiple
computational phases with inter-phase dependences (e.g.,
\texttt{correlation}: mean $\to$ normalization $\to$ correlation
matrix).  The analysis detects phases via cross-phase write conflicts
and recommends splitting into separate kernel launches connected by a
Python driver loop.  This is critical: LLMs frequently attempt to fuse
all phases into a single kernel, causing incorrect results or serial
execution.

\emph{Reduction patterns}: Reduction operations are mapped to Triton
primitives (\texttt{tl.sum()}, \texttt{tl.dot()}).  For kernels with
opposing reductions (e.g., \texttt{bicg}: row reduction on $A$ and
column reduction on $A^T$), the guidance recommends fusing both
reductions with row-major iteration for coalesced memory access.

\emph{Stencil synchronization}: For stencil computations, the analysis
detects cross-phase data dependences within a single CTA.  Triton's
vectorized stores go to L1 cache, and subsequent loads from overlapping
addresses may read stale data.  The guidance inserts
\texttt{tl.debug\_barrier()} calls between phases and includes a
pre-validation check that rejects generated code missing these barriers
before executing the expensive GPU test.

\paragraph{Generation and Refinement.}
The prompt is sent to Claude Sonnet~4~\cite{anthropic2025claude}, which
generates a Triton kernel implementation.  The generated code is tested
against a C reference compiled as a shared library (\texttt{-O0} to
preserve reference semantics): both implementations execute on
identical random inputs, and outputs are compared element-wise with
per-kernel tolerances that account for FP32 accumulation order
differences.

If validation fails, the system classifies the error and constructs a
targeted retry prompt.  For numerical mismatches, the feedback includes
the maximum absolute error and suggests index-bounds checking.  For
performance issues ($<$0.1$\times$ speedup), the feedback identifies
anti-patterns (serial grid, scalar loops, excessive launches).  Up to
10 attempts are made, with a context reset at attempt~6 to escape
fixation on broken approaches.

\sysname is implemented in approximately 14,000 lines of Python.
The analysis frontend comprises 18 modules built on \pet{} (which uses
isl~\cite{verdoolaege2010isl} for polyhedral operations) and LLVM~17
for fallback dependence analysis.


%% ===================================================================
\section{Evaluation}
\label{sec:eval}

All experiments run on an NVIDIA RTX~3090 GPU (24\,GB,
936\,GB/s) with a dual-socket Xeon Gold 5218R (80 cores,
${\sim}$140\,GB/s), PyTorch~2.5.1, Triton~3.1.0, CUDA~12.4.
The LLM is Claude Sonnet~4~\cite{anthropic2025claude}.
Polybench/application speedups are vs.\ single-threaded C
(\texttt{-O2}); TSVC additionally reports speedup vs.\ an
80-thread OpenMP baseline (\texttt{-fopenmp}).
All speedups are kernel-only with data pre-staged on GPU,
averaged over 10--50 iterations after warmup.

\subsection{Polybench/C (30 Kernels)}

\begin{table}[t]
  \caption{Polybench/C results.  WA = with analysis, NA = no analysis.}
  \label{tab:results_1x}
  \centering\small
  \begin{tabular}{lcc}
    \toprule
    \textbf{Metric} & \textbf{WA} & \textbf{NA} \\
    \midrule
    Pass rate              & 30/30 (100\%) & 28/30 (93\%) \\
    Median speedup         & 2.15$\times$  & 1.06$\times$ \\
    Mean speedup           & 3.44$\times$  & 1.52$\times$ \\
    Kernels $>1\times$     & 20/30         & 14/28         \\
    \midrule
    \multicolumn{3}{l}{\emph{Head-to-head (28 common):}} \\
    WA wins ($>$10\% faster) & \multicolumn{2}{c}{12} \\
    NA wins ($>$10\% faster) & \multicolumn{2}{c}{6 (3 race conditions)}  \\
    Ties (within 10\%)       & \multicolumn{2}{c}{10}  \\
    \bottomrule
  \end{tabular}
\end{table}

Table~\ref{tab:results_1x} summarizes results at standard
(1$\times$) problem sizes.  All 30 kernels pass with analysis
guidance; without it, two stencil kernels fail due to
parallelizing dimensions with carried dependences.
The 2.15$\times$ median speedup is limited by Polybench's
small defaults ($N$\,=\,60--120), where GPU launch overhead
dominates.  Compute-bound kernels show strong speedups
(\texttt{heat\_3d}: 12.36$\times$, \texttt{deriche}:
7.44$\times$, \texttt{covariance}: 7.40$\times$).

Of 6 apparent NA wins, 3 (\texttt{lu}, \texttt{ludcmp}, \texttt{bicg})
contain confirmed race conditions---the NA code uses
\texttt{grid=(N,)} for LU decomposition where row~$i$ reads from
row~$k$ ($k<i$) that may not yet be computed by its CTA, producing up
to $10^{33}$ variation across runs.  The remaining 3 NA wins reflect
nondeterministic LLM code quality, not structural differences.

\paragraph{Scaling.}

\begin{table}[t]
  \caption{Polybench scaling: 1$\times$ to 8$\times$ problem sizes
    ($N$: 60$\to$480, matrices: 120$\to$960).
    ``ST'' = single-threaded C; ``OMP'' = 80-thread OpenMP.}
  \label{tab:results_8x}
  \centering\small
  \begin{tabular}{lccc}
    \toprule
    \textbf{Metric} & \textbf{1$\times$} & \textbf{8$\times$ (ST)} & \textbf{8$\times$ (OMP)} \\
    \midrule
    Median speedup     & 2.15$\times$ & 28.97$\times$ & 7.98$\times$ \\
    Kernels $>1\times$ & 20/30        & 22/26         & 21/27        \\
    Max speedup        & 12.36$\times$ & 1826$\times$ & 1464$\times$ \\
    \bottomrule
  \end{tabular}
\end{table}

At 8$\times$ sizes (Table~\ref{tab:results_8x}), median speedup
rises to 28.97$\times$ over single-threaded C.  Against 80-thread
OpenMP, the GPU still achieves 7.98$\times$ median, confirming
genuine advantage beyond CPU parallelism.  Matrix-multiply kernels
reach 1464$\times$ (\texttt{3mm}) even vs.\ OpenMP.  Four kernels
time out (inherently sequential, \texttt{grid=(1,)}).

\paragraph{Guidance Impact.}

\begin{table}[t]
  \caption{Impact of analysis-triggered guidance categories.
    ``Before'' and ``After'' are speedups at 1$\times$ sizes
    except where noted.}
  \label{tab:guidance_fixes}
  \centering\small
  \begin{tabular}{llcc}
    \toprule
    \textbf{Guidance} & \textbf{Kernel} & \textbf{Before} & \textbf{After} \\
    \midrule
    Multi-phase split      & correlation & 0.19$\times$ & 2.50$\times$ \\
    WAR sequential         & nussinov    & FAIL          & PASS         \\
    Triangular grid        & lud$^\dagger$ & 0.14$\times$ & 13.18$\times$ \\
    Fused reductions       & bicg        & 1.04$\times$ & 3.51$\times$ \\
    Stencil barriers       & jacobi\_1d  & 0.13$\times$ & 1.78$\times$ \\
    IIR phase fusion       & deriche     & 0.28$\times$ & 7.19$\times$ \\
    Tridiag.\ sweep        & adi         & 0.24$\times$ & 2.27$\times$ \\
    \bottomrule
    \multicolumn{4}{l}{\scriptsize $^\dagger$Rodinia benchmark.}
  \end{tabular}
\end{table}

Table~\ref{tab:guidance_fixes} shows the impact of individual guidance
categories developed through iterative analysis of failure modes.
Each category is triggered by analysis properties (e.g., cross-phase
dependences, triangular access patterns), not per-kernel heuristics.
The largest improvements come from multi-phase decomposition
(correlation: 13$\times$ improvement) and the triangular grid
recommendation (lud: 94$\times$ improvement).  The IIR phase-fusion
guidance detects that the forward, backward, and combination phases
of a deriche filter operate on independent rows/columns, fusing all
three into a single kernel with \texttt{grid=(W,)} or
\texttt{grid=(H,)} (2 launches instead of 6).  The tridiagonal sweep
guidance identifies that ADI has one parallelizable spatial dimension
and one with sequential recurrence, recommending \texttt{grid=(N-2,)}
with one CTA per row.  At 8$\times$ sizes, a 2D stencil threshold
adjustment (multi-CTA grids when $\geq$2 valid dimensions exist)
improved \texttt{jacobi\_2d} to 36.82$\times$ and \texttt{fdtd\_2d}
to 15.43$\times$.

\paragraph{Limitations.}
Eleven kernels achieve $<$1$\times$ at standard sizes.  Five are
inherently sequential (\texttt{durbin}: 0.43$\times$,
\texttt{gramschmidt}: 0.35$\times$, \texttt{trisolv}: 0.58$\times$,
\texttt{ludcmp}: 0.36$\times$, \texttt{nussinov}: 0.06$\times$), where
loop-carried dependences on most dimensions limit parallelism; the
analysis correctly identifies these constraints and the LLM generates
working but slow GPU implementations.

Two kernels (\texttt{jacobi\_2d}: 0.71$\times$, \texttt{fdtd\_2d}:
0.34$\times$) have parallelizable spatial dimensions but at Polybench's
default sizes ($N$=20--30), launching 2--4 kernels per
timestep$\times$20--40 timesteps creates 80--160 launches that dominate
execution time.  At 8$\times$ sizes, these same kernels achieve
36.82$\times$ and 15.43$\times$, confirming that the parallelization
strategy is correct and the 1$\times$ regression is purely a
problem-size effect.  These $<$1$\times$ results are \emph{correct}
failures---preferable to the incorrect parallelization that the NA
baseline produces for \texttt{jacobi\_2d} and \texttt{seidel\_2d}.


\subsection{TSVC (151 Kernels)}
\label{sec:tsvc}

To evaluate breadth, we apply \sysname to the 151-kernel TSVC (Test
Suite for Vectorizing Compilers)~\cite{callahan1988tsvc} benchmark,
which covers linear traversals, reductions, indirect addressing,
control flow, and loop-carried dependences.  Unlike Polybench's
multi-dimensional stencils, TSVC kernels are predominantly 1D loops
operating on flat arrays, testing whether the analysis pipeline handles
simpler but diverse loop patterns.

\begin{table}[t]
  \caption{TSVC results at three data sizes.  ``ST'' = single-threaded
    C (\texttt{-O2}), ``OMP'' = OpenMP C (\texttt{-O2 -fopenmp},
    80~threads). Pass rate: 150/151 (99.3\%).}
  \label{tab:tsvc}
  \centering\small
  \begin{tabular}{lccc}
    \toprule
    & \textbf{32\,KB} & \textbf{32\,MB} & \textbf{3200\,MB} \\
    \midrule
    Per-array size     & 128\,KB & 32\,MB & 3.2\,GB \\
    Kernels tested     & 84      & 27     & 15      \\
    \midrule
    GPU/ST median      & 0.11$\times$ & 46.3$\times$ & 73.2$\times$ \\
    GPU/ST $>$1$\times$ & 1/84        & 25/27        & 14/15        \\
    \midrule
    GPU/OMP median     & 0.13$\times$ & 9.0$\times$  & 14.2$\times$ \\
    GPU/OMP $>$1$\times$& 1/84        & 24/27        & 14/15        \\
    \bottomrule
  \end{tabular}
\end{table}

Table~\ref{tab:tsvc} shows TSVC results across three data sizes.
\sysname passes 150 of 151 kernels (99.3\%), with 120 on the first
attempt.  At the default 32\,KB per array ($N$=32,000), GPU launch
overhead dominates and the median speedup is 0.11$\times$---neither
the GPU nor multi-threaded CPU can amortize overhead at this scale.

At 32\,MB ($N$=8M), the GPU achieves 46.3$\times$ median over
single-threaded C and 9.0$\times$ over 80-thread OpenMP.  At
3200\,MB ($N$=800M), the GPU advantage grows to 73.2$\times$ (ST)
and 14.2$\times$ (OMP), as the GPU's memory bandwidth (936\,GB/s on
RTX~3090) decisively outperforms the dual-socket Xeon's bandwidth
($\sim$140\,GB/s).  Serial kernels with loop-carried dependences are
excluded from the 32\,MB and 3200\,MB benchmarks since neither CPU
nor GPU can parallelize them.


\subsection{Application Kernels}
\label{sec:realworld}

\begin{table}[t]
  \caption{Application kernel results.  ``---'' = NA failed all attempts.}
  \label{tab:realworld}
  \centering\small
  \begin{tabular}{llcc}
    \toprule
    \textbf{Source} & \textbf{Kernel} & \textbf{WA} & \textbf{NA} \\
    \midrule
    Rodinia & SRAD       & 24.86$\times$ & 3.50$\times$ \\
    Rodinia & lud        & 13.18$\times$ & 15.52$\times$ \\
    Rodinia & gaussian   & 3.35$\times$  & 3.61$\times$ \\
    Rodinia & lavaMD     & 1.47$\times$  & 0.56$\times$ \\
    Rodinia & hotspot    & 1.05$\times$  & 1.22$\times$ \\
    Rodinia & pathfinder & 0.11$\times$  & --- \\
    \midrule
    miniMD  & LJ force   & 25.82$\times$ & 23.00$\times$ \\
    miniFE  & SpMV       & 8.95$\times$  & 7.27$\times$ \\
    \bottomrule
  \end{tabular}
\end{table}

To evaluate generalization beyond synthetic benchmarks, we apply
\sysname to 6 Rodinia kernels~\cite{che2009rodinia} and 2 ECP proxy
application kernels~\cite{heroux2009mantevo}---CSR sparse matrix-vector
multiply (SpMV) from miniFE and Lennard-Jones force computation from
miniMD (Table~\ref{tab:realworld}).  All 8 pass correctness with
analysis guidance (6 on the first attempt).

SRAD (24.9$\times$) has three computational phases per iteration;
the analysis detects cross-phase dependences and recommends splitting
into separate kernels with parallel grids, yielding 7.1$\times$
improvement over the NA baseline (3.5$\times$).
The LJ force kernel from miniMD computes pairwise forces using a full
neighbor list; the outer loop over atoms is embarrassingly parallel.
With analysis guidance the LLM iterates only actual neighbors
(25.8$\times$), while the NA version wastes cycles on padded entries
(23.0$\times$).
lavaMD improves from 0.56$\times$ (NA) to 1.47$\times$ (WA) as
analysis identifies the parallelizable particle dimension, while the
NA baseline fails with type-mismatch errors on 4 of 5 attempts.
SpMV achieves 8.95$\times$ with analysis vs.\ 7.27$\times$ without,
as guidance helps select better block sizes for the row-parallel pattern.
For lud and hotspot, both versions produce near-identical Triton
times; the speedup differences reflect C reference timing variance.


%% ===================================================================
\section{Related Work}
\label{sec:related}

\paragraph{LLM Code Generation for GPUs.}
Chen et al.~\cite{chen2024chatgpt} evaluate ChatGPT on CUDA kernel
generation, finding frequent correctness issues including data races and
incorrect memory management.
KernelBench~\cite{ouyang2024kernelbench} benchmarks LLM-generated
Triton/CUDA on 250 tasks, reporting 20--40\% functional correctness
without guidance.  Our work differs fundamentally: rather than relying
on the LLM's implicit knowledge of parallelism, we provide explicit
compiler-derived constraints, achieving 100\% correctness on
Polybench/C and 99\% on TSVC.

\paragraph{Polyhedral Compilation and Auto-Scheduling.}
Polyhedral frameworks such as Pluto~\cite{bondhugula2008pluto} and
PPCG~\cite{verdoolaege2013ppcg} automatically generate parallel code
from affine loop nests.  PPCG targets CUDA directly but is limited to
affine programs and produces low-level code that lacks Triton's
block-level optimizations (automatic memory coalescing, shared memory
tiling).  TVM~\cite{chen2018tvm} and Halide~\cite{ragan2013halide}
use auto-scheduling with learned cost models for tensor operations.
\sysname is complementary: rather than replacing the compiler,
it uses polyhedral analysis to extract \emph{what} parallelism is
safe, then delegates the \emph{how} to the LLM, which can handle
non-affine constructs, irregular access patterns, and high-level
Triton abstractions that polyhedral compilers cannot target directly.


%% ===================================================================
\section{Conclusion}
\label{sec:conclusion}

We presented \sysname, a system that bridges polyhedral compiler
analysis with LLM-based GPU kernel generation.  By translating
parallelization constraints, dependence information, and memory access
patterns into structured prompt guidance, \sysname enables an LLM to
generate correct and efficient Triton kernels for all 30 Polybench/C
benchmarks (2.15$\times$ median), 150 of 151 TSVC kernels (9$\times$
median over 80-thread OpenMP at realistic data sizes), and 8
application kernels with up to 25.8$\times$ speedup.

The ablation study reveals two key findings: (1)~analysis guidance
prevents correctness failures that stem from unsafe parallelization
of dependent loop dimensions, and (2)~unguided code frequently contains
latent race conditions---producing results that pass tolerance checks
by chance but exhibit nondeterministic errors at larger scales.
Three of six apparent NA ``wins'' in Polybench were confirmed to
contain such race conditions via repeated execution.

For application kernels, analysis guidance is essential for
complex multi-phase computations (SRAD: 7$\times$ improvement
over unguided) and beneficial even for naturally parallel patterns
(LJ~force: guidance promotes efficient neighbor-list traversal).
No application kernel shows a Triton-time regression with analysis
guidance enabled.

Future work includes extending the analysis to non-affine programs,
supporting multi-GPU parallelization, and integrating autotuning
for block size selection.

%% ===================================================================
%% References
\bibliographystyle{ACM-Reference-Format}
\begin{thebibliography}{99}

\bibitem{tillet2019triton}
P.~Tillet, H.~T. Kung, and D.~Cox.
\newblock Triton: An intermediate language and compiler for tiled neural
  network computations.
\newblock In \emph{Proc. MAPL}, 2019.

\bibitem{chen2024chatgpt}
Y.~Chen, S.~Huang, D.~Zheng, et~al.
\newblock {ChatGPT} for {CUDA} programming: Evaluating {LLM}-generated {GPU}
  code.
\newblock \emph{arXiv preprint arXiv:2404.12000}, 2024.

\bibitem{ouyang2024kernelbench}
A.~Ouyang, S.~Zheng, et~al.
\newblock {KernelBench}: Can {LLMs} write efficient {GPU} kernels?
\newblock \emph{arXiv preprint arXiv:2404.12489}, 2024.

\bibitem{verdoolaege2012polyhedral}
S.~Verdoolaege and T.~Grosser.
\newblock Polyhedral extraction tool.
\newblock In \emph{Proc. IMPACT}, 2012.

\bibitem{llvm}
C.~Lattner and V.~Adve.
\newblock {LLVM}: A compilation framework for lifelong program analysis \&
  transformation.
\newblock In \emph{Proc. CGO}, 2004.

\bibitem{polybench}
L.-N. Pouchet.
\newblock Polybench/C: The polyhedral benchmark suite, version 4.2.1.
\newblock \url{https://web.cs.ucla.edu/~pouchet/software/polybench/}, 2016.

\bibitem{anthropic2025claude}
Anthropic.
\newblock Claude {Sonnet} 4 model card.
\newblock \url{https://www.anthropic.com/claude}, 2025.

\bibitem{verdoolaege2010isl}
S.~Verdoolaege.
\newblock isl: An integer set library for the polyhedral model.
\newblock In \emph{Proc. ICMS}, 2010.

\bibitem{chen2018tvm}
T.~Chen, T.~Moreau, Z.~Jiang, et~al.
\newblock {TVM}: An automated end-to-end optimizing compiler for deep
  learning.
\newblock In \emph{Proc. OSDI}, 2018.

\bibitem{ragan2013halide}
J.~Ragan-Kelley, C.~Barnes, A.~Adams, et~al.
\newblock Halide: A language and compiler for optimizing parallelism,
  locality, and recomputation in image processing pipelines.
\newblock In \emph{Proc. PLDI}, 2013.

\bibitem{bondhugula2008pluto}
U.~Bondhugula, A.~Hartono, J.~Ramanujam, and P.~Sadayappan.
\newblock A practical automatic polyhedral parallelizer and locality
  optimizer.
\newblock In \emph{Proc. PLDI}, 2008.

\bibitem{verdoolaege2013ppcg}
S.~Verdoolaege, J.~Carlos~Juega, A.~Cohen, et~al.
\newblock Polyhedral parallel code generation for {CUDA}.
\newblock \emph{ACM TACO}, 9(4), 2013.

\bibitem{che2009rodinia}
S.~Che, M.~Boyer, J.~Meng, et~al.
\newblock Rodinia: A benchmark suite for heterogeneous computing.
\newblock In \emph{Proc. IISWC}, 2009.

\bibitem{heroux2009mantevo}
M.~A.~Heroux, D.~W.~Doerfler, P.~S.~Crozier, et~al.
\newblock Improving performance via mini-applications.
\newblock Sandia National Laboratories Technical Report SAND2009-5574, 2009.

\bibitem{callahan1988tsvc}
D.~Callahan, J.~Dongarra, and D.~Levine.
\newblock Vectorizing compilers: A test suite and results.
\newblock In \emph{Proc. Supercomputing}, 1988.

\end{thebibliography}

\end{document}
